\documentclass[10pt,a4paper]{amsart}
\usepackage[margin=19mm]{geometry}
\usepackage{amsmath,amssymb,mathtools}
\usepackage[scr=rsfs,cal=euler]{mathalfa}
\usepackage{xfrac}
\usepackage[unicode,hidelinks,bookmarks=false]{hyperref}
\newcommand{\ie}{i.e.\ }
\newcommand{\cf}{cf.\ }
\newcommand{\resp}{respectively\ }
\newcommand{\Subsection}{\subsection}
\providecommand{\nobreakdash}{\nobreak\hbox{-}}
\setlength{\emergencystretch}{2em}
\multlinegap0pt
\usepackage{enumitem}
\newtheorem{theorem}{Theorem}
\newtheorem{lemma}[theorem]{Lemma}
\newtheorem*{bhfirst}{Theorem A}
\title[$C^{1,\alpha}$ regularity of the free boundaries for the American chooser option (Theorem 4.9)]{The Obstacle Problem Arising from the American Chooser Option: $C^{1,\alpha}$ regularity of the two free boundaries (Theorem 4.9)}
\author{Gugyum Ha, Junkee Jeon, Jihoon Ok}
\date{Source: arXiv:2506.03623v4 (CC BY 4.0)}
\begin{document}
\begin{abstract}
We study the finite-horizon obstacle problem associated with the American
chooser option. Its obstacle is the maximum of two independently determined
American call and put values and has a spatial kink where these values
cross. Using truncation, penalization, and comparison, we construct a local
strong Sobolev solution and prove uniqueness within an exponential-growth
class. We show that contact occurs only on the smooth exercise-payoff
branches and that both contact regions have nonempty sections at every
time. We then establish strict monotonicity, terminal limits, and local
Lipschitz continuity of the two free boundaries. Localized one-phase
equations, quantitative nondegeneracy, and parabolic boundary Harnack
estimates yield $C^\infty$ regularity on the open time interval.


\end{abstract}
\maketitle

\noindent\emph{Source and licence.} The Obstacle Problem Arising from the American Chooser Option. Version arXiv:2506.03623v4 (CC BY 4.0), distributed under the Creative Commons Attribution 4.0 International licence, as recorded in the arXiv licence metadata for this exact version and shown on the abstract page https://arxiv.org/abs/2506.03623v4. Full rights notice: licenses/arxiv-2506.03623-CC-BY-4.0.txt.\par\smallskip

\noindent\textit{Editorial scope.} Counted unit: the theorem printed as Theorem 4.9 of the article (source0.tex lines 2996--3002, source label \texttt{C1alpha free boundary}), delivered together with the article's own complete original proof of it (source0.tex lines 3004--3277, one \texttt{proof} environment of 274 lines that proves the local $C^{1,\alpha}$ regularity of both free boundaries with no gap and no deferral). No mathematics is written, repaired or re-derived: statement and proof are the article's own text line for line, in the article's own notation ($V$, $x_p^{\rm ch}$, $x_c^{\rm ch}$, $x_{p,n}$, $\mathcal C$, $\mathcal E_C^{\rm ch}$, $\Omega_T$, $\mathcal L$, $Q_r(\tau_0,x_0)$). Required context retained uncounted: the article's abstract (lines 184--194); the identity (Vtau boundary) recording that the time derivative of $V$ vanishes at both boundaries (lines 2441--2447); the article's Lemma 4.6, the uniform estimate $0\leq\partial_\tau V\leq C_I(\partial_xV+e^x)$ on compact time intervals, in statement form (lines 2471--2497); the display (put level curve) defining the level curves $x_{p,n}$ by $V(\tau,x_{p,n}(\tau))-(K_p-e^{x_{p,n}(\tau)})=1/n$ (lines 2666--2670); the article's Lemma 4.8, the quantitative nondegeneracy estimate at the put boundary, in statement form (lines 2807--2823); and the article's Theorem A, the parabolic boundary Harnack inequality with its attribution to Kukuljan, in statement form (lines 2975--2984), which the counted proof invokes by name three times. Every definition, numbered display and label of those lines is retained unchanged. Omitted from this package and not redistributed: the article's preamble, title page, subject classification and keywords (lines 1--183); the end of its front matter (lines 195--199); its whole introduction (lines 200--346); its second section with the notations, the financial background of the American chooser option, the change of variables and the properties of American call and put options (lines 347--832); its third section with the truncation and penalization construction, the existence and uniqueness of the strong solution and the estimate for the derivatives of $V$ (lines 833--1798); the parametrization and the properties of the free boundary in the fourth section (lines 1800--2440); and, inside the fourth section's regularity subsection, everything that is not one of the retained excerpts, namely the derivation of the level curves and of their Lipschitz bound (lines 2448--2470, 2671--2806), the intervening monotonicity and limit results (lines 2498--2665, 2824--2974) and the block of named theorems quoted from the literature that precedes the counted theorem (lines 2985--2995); the whole appendix with the formulation of the model, the verification, the proof of the call-put inequality and the conditions on the operator (lines 3278--3924); the article's bibliography (lines 3925--4026); and its closing end-document line (line 4027). Nothing inside a retained excerpt is omitted or altered: every retained body is delivered exactly as the article's lines are written, so no omission receipt of any kind accompanies this package. Disclosed dependence on omitted material: the retained displays and the counted proof use the level-curve indices $x_{p,n}$ and the continuation and exercise sets of the article's fourth section, whose definitions live in the omitted part of that section, and the counted proof also speaks of the local estimates proved in the omitted earlier subsections; the two exact statements that the proof actually cites (Lemma 4.6 and Lemma 4.8) and the two identities it uses are all retained. Citations: the retained excerpts carry exactly two citation commands, one in the optional argument of the retained Theorem A and one inside the counted proof, and both cite the same key of the article's own bibliography, \texttt{Kukuljan2022}. That entry is transcribed verbatim from the article's own in-source bibliography (lines 3925--4026, 29 entries) with the article's own printed label, so the delivered citation marks read [16] as the article's do. No other entry of the article's bibliography is transcribed or redistributed. Reference adaptation: none was needed and none was made. Every reference inside the delivered unit resolves inside it: the labels FB nondeg, derivative estimate FB, Vtau boundary and put level curve of the retained statements and displays, all cited by the counted proof, together with the label C1alpha free boundary of the counted statement itself. No unresolved reference or citation is produced. Printed slips kept literal: no printed word, symbol, hypothesis or number of the counted statement or of its proof was altered; in particular the retained text keeps the article's own use of the old-style font switch $\rm$ in $x_p^{\rm ch}$ and $x_c^{\rm ch}$. Layout and class: the article's own class is \texttt{amsart} with the options 11pt, a4paper and leqno, and it loads among others inputenc, fontenc, times, amsmath, amssymb, amsthm, enumitem, graphicx and hyperref; no publisher class or style file is used or shipped, so none travels with this package and none is redistributed. The wrapper is the lane's pinned \texttt{amsart} editorial wrapper at 10pt on A4, which declares the theorem-like environments the retained text needs (\texttt{theorem}, \texttt{lemma} sharing the theorem counter, and the unnumbered \texttt{bhfirst} for the article's Theorem A), loads the \texttt{enumitem} package that the counted proof's roman-numbered list needs, and needs no macro from the article's preamble because every control sequence in the retained spans is standard. The delivered numbering is the unit's own: the retained Lemma 4.6 prints as Lemma 1, the retained Lemma 4.8 as Lemma 2, the retained Theorem A stays unnumbered under its own name, and the counted Theorem 4.9 prints as Theorem 3. The article's leqno option (equation numbers on the left) and its Times text font are not reproduced by the wrapper, which numbers equations on the right in Latin Modern; those are page-appearance differences of the wrapper only, and no formula, symbol, hypothesis or argument changes. Assets: no retained span contains a figure environment, a graphics-inclusion command, a PDF-page inclusion, a table or a TikZ picture, and the article's own figures are confined to its omitted introduction and numerical discussion; no image file is copied into this package, the package declares no asset dependency, and no image file is redistributed with it. Licence: version arXiv:2506.03623v4 of this article is distributed under the Creative Commons Attribution 4.0 International licence, as recorded in the arXiv licence metadata for this exact version and shown on the abstract page https://arxiv.org/abs/2506.03623v4. A full rights notice is delivered at licenses/arxiv-2506.03623-CC-BY-4.0.txt. Characters: every retained excerpt is pure ASCII, so the delivered engine, XeLaTeX with its default Latin Modern text font, reports no missing character. No glyph is substituted and no word is rewritten.
\begin{equation}\label{Vtau boundary}
    \partial_\tau V(\tau,x_p^{\rm ch}(\tau))
    =
    \partial_\tau V(\tau,x_c^{\rm ch}(\tau))
    =0,
    \qquad \tau\in(0,T).
\end{equation}

\begin{lemma}\label{derivative estimate FB}
For every compact interval $I\Subset(0,T)$, there exists a constant
$C_I>0$ such that
\begin{equation}\label{put derivative estimate}
    0\leq\partial_\tau V(\tau,x)
    \leq
    C_I\bigl(\partial_xV(\tau,x)+e^x\bigr)
\end{equation}
for
\[
    \tau\in I,
    \qquad
    x_p^{\rm ch}(\tau)<x\leq\ln K_p,
\]
and
\begin{equation}\label{call derivative estimate}
    0\leq\partial_\tau V(\tau,x)
    \leq
    C_I\bigl(e^x-\partial_xV(\tau,x)\bigr)
\end{equation}
for
\[
    \tau\in I,
    \qquad
    \ln K_c\leq x<x_c^{\rm ch}(\tau).
\]
\end{lemma}

\begin{equation}\label{put level curve}
    V(\tau,x_{p,n}(\tau))
    -(K_p-e^{x_{p,n}(\tau)})
    =\frac1n,
\end{equation}

\begin{lemma}\label{FB nondeg}
Let $I\Subset(0,T)$ be a compact interval.
Then there exist constants $\delta>0$ and $c>0$ such that
\begin{align}
    \partial_xV(\tau,x)+e^x
    &\geq
    c\bigl(x-x_p^{\rm ch}(\tau)\bigr),
    &&x_p^{\rm ch}(\tau)<x<x_p^{\rm ch}(\tau)+\delta,
    \label{put nondeg}\\
    e^x-\partial_xV(\tau,x)
    &\geq
    c\bigl(x_c^{\rm ch}(\tau)-x\bigr),
    &&x_c^{\rm ch}(\tau)-\delta<x<x_c^{\rm ch}(\tau),
    \label{call nondeg}
\end{align}
for all $\tau\in I$.
\end{lemma}

\begin{bhfirst}[{\cite[Theorem~1.2]{Kukuljan2022}}]
Suppose $D$ is a parabolic $C_p^1$ graph domain, $a$ is continuous
and uniformly positive and bounded, and $b$ is bounded.
Let bounded $u_i$ satisfy $\mathcal P_0u_i=f_i$ in $D$ and
$u_i=0$ on its local boundary, with bounded $f_i$.
If $|u_2|\geq c_2d$ and
$\|u_2\|_\infty+\|f_2\|_\infty\leq C_2$, then
$u_1/u_2\in C_p^{1-\varepsilon}$ up to the boundary of every smaller
cylinder, for each $0<\varepsilon<1$.
\end{bhfirst}

\begin{theorem}\label{C1alpha free boundary}
For every $\alpha\in(0,\frac12)$,
\[
    x_p^{\rm ch},\,x_c^{\rm ch}
    \in C_{\rm loc}^{1,\alpha}((0,T)).
\]
\end{theorem}

\begin{proof}
We first consider $x_p^{\rm ch}$.
Fix $\tau_0\in(0,T)$ and set
\[
    x_0:=x_p^{\rm ch}(\tau_0).
\]
Choose a compact interval $I\Subset(0,T)$ containing $\tau_0$ in its
interior. By Lemma~\ref{FB nondeg}, there exist $\delta>0$ and $c>0$
such that
\[
    \partial_xV(\tau,x)+e^x
    \geq
    c\bigl(x-x_p^{\rm ch}(\tau)\bigr)
\]
for $\tau\in I$ and
\[
    x_p^{\rm ch}(\tau)
    <x<
    x_p^{\rm ch}(\tau)+\delta.
\]

Let $L$ be a Lipschitz constant of $x_p^{\rm ch}$ on $I$, and define
the parabolic cylinder
\[
    Q_r(\tau_0,x_0)
    :=
    (\tau_0-r^2,\tau_0+r^2)
    \times(x_0-r,x_0+r).
\]
Choose $r>0$ sufficiently small so that
\[
    [\tau_0-4r^2,\tau_0+4r^2]\subset I,
    \qquad
    2r+4Lr^2<\delta,
\]
and
\[
    Q_{2r}(\tau_0,x_0)\cap\mathcal E_C^{\rm ch}
    =\emptyset.
\]
Then
\[
    Q_{2r}(\tau_0,x_0)\cap\mathcal C
    =
    Q_{2r}(\tau_0,x_0)
    \cap
    \{x>x_p^{\rm ch}(\tau)\}.
\]
Moreover, for every $(\tau,x)$ in this set,
\[
\begin{aligned}
    0
    &<
    x-x_p^{\rm ch}(\tau) \\
    &\leq
    |x-x_0|
    +
    |x_p^{\rm ch}(\tau)-x_p^{\rm ch}(\tau_0)| \\
    &<
    2r+4Lr^2
    <\delta.
\end{aligned}
\]
Consequently,
\[
    \partial_xV(\tau,x)+e^x
    \geq
    c\bigl(x-x_p^{\rm ch}(\tau)\bigr)
    \quad\text{in }
    Q_{2r}(\tau_0,x_0)\cap\mathcal C.
\]

In $Q_{2r}(\tau_0,x_0)\cap\mathcal C$, we have
\[
    (\partial_\tau-\mathcal L)(\partial_\tau V)=0,
    \qquad
    (\partial_\tau-\mathcal L)(\partial_xV+e^x)=qe^x.
\]
By the continuity of $\partial_\tau V$ and $\partial_xV$ up to
$x=x_p^{\rm ch}(\tau)$ from the continuation region, we have
\[
    \partial_\tau V=0,
    \qquad
    \partial_xV+e^x=0
    \quad\text{on }
    Q_{2r}(\tau_0,x_0)
    \cap\{x=x_p^{\rm ch}(\tau)\}.
\]

{
We check the hypotheses of Theorem~A one by one on a smaller cylinder.
\begin{enumerate}[label=(\roman*)]
\item \emph{Domain.} For $g(\tau)=x_p^{\rm ch}(\tau)$,
\[
 \frac{|g(\tau)-g(s)|}{|\tau-s|^{1/2}}
 \leq L|\tau-s|^{1/2}\longrightarrow0.
\]
There are no tangential spatial variables in one dimension. Thus a
time-Lipschitz graph is $C_p^1$ in the sense of
\cite[Definition~1.1]{Kukuljan2022}; in fact it belongs to
$C_p^\beta$ for every $1<\beta<2$.
\item \emph{Operator.} Put $a=\sigma^2/2$,
$b=r-q-\sigma^2/2$ and use
$u_1=e^{r\tau}V_\tau$ and
$u_2=e^{r\tau}(V_x+e^x)$. Then
$\mathcal P_0=\partial_\tau-a\partial_{xx}-b\partial_x$
has constant uniformly parabolic coefficients and no zero-order term.
\item \emph{Equations and bounds.} The preceding equations become
\[
 \mathcal P_0u_1=0,\qquad
 \mathcal P_0u_2=q e^{r\tau+x}.
\]
The sources are smooth and bounded locally. The local continuity of
$V_\tau$ and $V_x$ gives the required bounds for $u_1,u_2$.
\item \emph{Traces.} Equation~\eqref{Vtau boundary} and spatial
smooth fit give $u_1=u_2=0$ on the local boundary graph.
\item \emph{Nondegeneracy.} The same-time boundary point gives
$d(\tau,x)\leq x-g(\tau)$. Lemma~\ref{FB nondeg} therefore yields
$u_2\geq c_2d$, since $e^{r\tau}$ has a positive lower bound locally.
% For completeness, the reverse comparison of these distances also
% holds on smaller cylinders: for any nearby boundary point $(s,g(s))$,
% \[
%  x-g(\tau)\leq |x-g(s)|+L|\tau-s|
%  \leq C\bigl(|x-g(s)|+|\tau-s|^{1/2}\bigr).
% \]
% Taking the infimum gives $x-g(\tau)\leq Cd$.
\end{enumerate}
Theorem~A now applies, and its quotient $u_1/u_2$ is exactly
$V_\tau/(V_x+e^x)$. Taking $\varepsilon=1-2\alpha$ gives the
following estimate.
}
For every $\alpha\in(0,\frac12)$, setting
\[
    F(\tau,x)
    :=
    \frac{\partial_\tau V(\tau,x)}
    {\partial_xV(\tau,x)+e^x},
\]
we obtain
\[
    |F(\tau,x)-F(s,y)|
    \leq
    C\bigl(|\tau-s|^\alpha+|x-y|^{2\alpha}\bigr)
\]
for $(\tau,x),(s,y)\in Q_r(\tau_0,x_0)\cap\mathcal C$.

Recall from \eqref{put level curve} that
\[
    x_{p,n}'(\tau)
    =
    -F(\tau,x_{p,n}(\tau)).
\]
Since $x_p^{\rm ch}$ is Lipschitz continuous on $I$ with constant $L$,
by taking $r>0$ sufficiently small if necessary, we have
\[
    x_p^{\rm ch}(\tau)
    \leq x_0+Lr^2
    <x_0+r
    \qquad\text{for }|\tau-\tau_0|\leq r^2.
\]
Hence, the right lateral boundary of
$Q_r(\tau_0,x_0)$ lies strictly in the continuation region. Therefore,
\[
    m_r:=
    \min_{|\tau-\tau_0|\leq r^2}
    \left[
        V(\tau,x_0+r)-\bigl(K_p-e^{x_0+r}\bigr)
    \right]>0.
\]
Choosing $n$ sufficiently large so that $1/n<m_r$, the level curve
$x_{p,n}$ lies in $Q_r(\tau_0,x_0)\cap\mathcal C$ for
$|\tau-\tau_0|\leq r^2$.

Moreover, combining Lemma~\ref{derivative estimate FB} with the
level-curve identity above gives a uniform Lipschitz bound for
$x_{p,n}$. Therefore,
\[
    |x_{p,n}(\tau)-x_{p,n}(s)|
    \leq C|\tau-s|,
\]
where $C>0$ is independent of $n$. Using the H\"older estimate for
$F$, we then obtain
\[
\begin{aligned}
    |x_{p,n}'(\tau)-x_{p,n}'(s)|
    &=
    \left|
        F(\tau,x_{p,n}(\tau))
        -F(s,x_{p,n}(s))
    \right| \\
    &\leq
    C\left(
        |\tau-s|^\alpha
        +|x_{p,n}(\tau)-x_{p,n}(s)|^{2\alpha}
    \right) \\
    &\leq
    C\left(
        |\tau-s|^\alpha+|\tau-s|^{2\alpha}
    \right) \\
    &\leq
    C|\tau-s|^\alpha.
\end{aligned}
\]
Since the uniform Lipschitz bound also gives
$\|x_{p,n}'\|_{L^\infty}\leq C$, it follows that
\[
    \|x_{p,n}'\|_{C^\alpha}\leq C.
\]
As $x_{p,n}$ lies in $Q_r(\tau_0,x_0)$, we consequently have
\[
    \|x_{p,n}\|_{C^{1,\alpha}}\leq C
\]
on $[\tau_0-r^2,\tau_0+r^2]$, with $C$ independent of $n$.

By the Arzel\`a--Ascoli theorem, a subsequence converges in $C^1$.
Since $x_{p,n}(\tau)\to x_p^{\rm ch}(\tau)$ for each $\tau$, the limit
is $x_p^{\rm ch}$. Passing to the limit in the preceding estimate
therefore gives
\[
    x_p^{\rm ch}\in C^{1,\alpha}
    \bigl((\tau_0-r^2,\tau_0+r^2)\bigr).
\]
Since $\tau_0\in(0,T)$ was arbitrary,
\[
    x_p^{\rm ch}\in C_{\rm loc}^{1,\alpha}((0,T)).
\]

We next consider $x_c^{\rm ch}$.
Fix $\tau_0\in(0,T)$ and set
\[
    x_0:=x_c^{\rm ch}(\tau_0).
\]
By Lemma~\ref{FB nondeg} and the local Lipschitz continuity of
$x_c^{\rm ch}$, we may choose $r>0$ sufficiently small so that
\[
    Q_{2r}(\tau_0,x_0)\cap\mathcal C
    =
    Q_{2r}(\tau_0,x_0)\cap\{x<x_c^{\rm ch}(\tau)\},
\]
and
\[
    e^x-\partial_xV(\tau,x)
    \geq
    c\bigl(x_c^{\rm ch}(\tau)-x\bigr)
    \quad\text{in }
    Q_{2r}(\tau_0,x_0)\cap\mathcal C.
\]
In this set,
\[
    (\partial_\tau-\mathcal L)(\partial_\tau V)=0,
    \qquad
    (\partial_\tau-\mathcal L)(e^x-\partial_xV)=qe^x,
\]
with
\[
    \partial_\tau V=0,
    \qquad
    e^x-\partial_xV=0
    \quad\text{on }
    Q_{2r}(\tau_0,x_0)\cap\{x=x_c^{\rm ch}(\tau)\}.
\]
{
For this side use $u_1=e^{r\tau}V_\tau$ and
$u_2=e^{r\tau}(e^x-V_x)$. The same operator, source, and trace
checks apply, while $d\leq x_c^{\rm ch}(\tau)-x$ gives the
required lower bound. Reflection of the spatial coordinate puts the
domain in the graph orientation of Theorem~A.
} Applying Theorem~A and repeating the preceding argument with the
level curves $x_{c,n}$, we obtain
\[
    x_c^{\rm ch}\in C_{\rm loc}^{1,\alpha}((0,T)).
\]
Since $\alpha\in(0,\frac12)$ was arbitrary, the proof is complete.
\end{proof}

\begin{thebibliography}{99}
\bibitem[16]{Kukuljan2022} T.~Kukuljan, Higher order parabolic boundary Harnack inequality in $C^1$ and $C^{k,\alpha}$ domains, \textit{Discrete and Continuous Dynamical Systems}, Volume 42, Number 6, (2022), 2667--2698.
\end{thebibliography}
\end{document}
